\documentclass[12pt,a4paper]{article}
\usepackage[margin=2.5cm]{geometry}
\usepackage{amsmath,amssymb,amsfonts}
\usepackage{physics}
\usepackage{enumitem}
\usepackage{fancyhdr}
\usepackage{lastpage}
\usepackage{graphicx}
\usepackage{multicol}
\usepackage{array}
\usepackage{tabularx}
\usepackage{tikz}

% Page style
\pagestyle{fancy}
\fancyhf{}
\renewcommand{\headrulewidth}{0.4pt}
\renewcommand{\footrulewidth}{0.4pt}
\lhead{PHYS10012}
\rhead{Mock Examination 4}
\cfoot{Page \thepage\ of \pageref{LastPage}}

% Question numbering
\newcounter{questioncounter}
\newcommand{\question}[1]{%
  \stepcounter{questioncounter}%
  \vspace{0.5cm}%
  \noindent\textbf{Question \thequestioncounter.} \hfill [#1 marks]%
  \vspace{0.2cm}%
  \par%
}

% Sub-part environments
\newlist{parts}{enumerate}{2}
\setlist[parts,1]{label=(\alph*),leftmargin=2em,itemsep=0.3cm}
\setlist[parts,2]{label=(\roman*),leftmargin=2em,itemsep=0.2cm}

% Mark allocation in sub-parts
\renewcommand{\marks}[1]{\hfill [#1]}

% MCQ environment
\newcounter{mcqcounter}
\newcommand{\mcq}[1]{%
  \stepcounter{mcqcounter}%
  \vspace{0.3cm}%
  \noindent\textbf{\themcqcounter.} #1 \hfill [2 marks]%
  \vspace{0.1cm}%
}
\newlist{choices}{enumerate}{1}
\setlist[choices]{label=\Alph*.,leftmargin=2.5em,itemsep=0.1cm}

% Dirac notation shortcuts
\renewcommand{\ket}[1]{\left|#1\right\rangle}
\renewcommand{\bra}[1]{\left\langle #1\right|}
\renewcommand{\braket}[2]{\left\langle #1\middle|#2\right\rangle}
\renewcommand{\ketbra}[2]{\left|#1\right\rangle\!\left\langle #2\right|}
\renewcommand{\expect}[1]{\left\langle #1\right\rangle}

\begin{document}

% Title block
\begin{center}
{\Large\textbf{UNIVERSITY OF BRISTOL}}\\[0.3cm]
{\large Faculty of Science}\\[0.5cm]
{\Large\textbf{MOCK EXAMINATION 4}}\\[0.3cm]
{\large\textbf{Core Physics I --- Classical, Quantum and Thermal Physics}}\\[0.2cm]
{\large Module Code: PHYS10012}\\[0.5cm]
{\large Duration: 3 Hours}\\[0.3cm]
{\large Total Marks: 100}\\[0.5cm]
\end{center}

\noindent\rule{\textwidth}{0.4pt}

\vspace{0.3cm}
\noindent\textbf{Instructions:}
\begin{itemize}[itemsep=0.1cm]
\item Answer \textbf{ALL} questions.
\item \textbf{Section A} (Oscillations and Waves): 10 multiple-choice questions, 20 marks total.
\item \textbf{Section B} (Quantum Physics): 10 multiple-choice questions, 20 marks total.
\item \textbf{Section C}: 4 long-answer questions, 60 marks total.
\item An approved calculator is permitted.
\item A formula sheet is provided at the end of this paper.
\item Show all working for Section C. Marks will be awarded for clear, logical solutions.
\end{itemize}

\noindent\rule{\textwidth}{0.4pt}
\newpage

% ============================================================
% SECTION A: OSCILLATIONS AND WAVES MCQ
% ============================================================
\section*{Section A: Oscillations and Waves}
\noindent\textit{Answer all 10 questions. Each question is worth 2 marks.}\\[0.3cm]

\setcounter{mcqcounter}{0}

% QUESTION A1: SHM pendulum period
\mcq{A simple pendulum of length $L$ has period $T$. What length of pendulum gives a period of $2T$?}
\begin{choices}
\item $2L$
\item $4L$
\item $\sqrt{2}\,L$
\item $L/2$
\end{choices}

% QUESTION A2: Energy in SHM
\mcq{A simple harmonic oscillator has spring constant $k = 200$~N/m and amplitude $A = 0.1$~m. What is the total energy of the oscillator?}
\begin{choices}
\item $0.5$~J
\item $1.0$~J
\item $2.0$~J
\item $10$~J
\end{choices}

% QUESTION A3: Critical damping
\mcq{A damped harmonic oscillator has natural angular frequency $\omega_0$ and damping coefficient $\gamma = b/m$. Critical damping occurs when:}
\begin{choices}
\item $\gamma = \omega_0$
\item $\gamma = 2\omega_0$
\item $\gamma = \omega_0/2$
\item $\gamma = \omega_0^2$
\end{choices}

% QUESTION A4: Power absorption in forced oscillator
\mcq{A forced harmonic oscillator has natural frequency $\omega_0$ and is driven at angular frequency $\omega$. The average power absorbed from the driving force is maximum when:}
\begin{choices}
\item $\omega = \omega_0/\sqrt{2}$
\item $\omega = \omega_0\sqrt{1 - \gamma^2/(2\omega_0^2)}$
\item $\omega = \omega_0$
\item $\omega = 2\omega_0$
\end{choices}

% QUESTION A5: Wave speed on string
\mcq{A transverse wave travels along a string with speed $v$. If the tension in the string is quadrupled while the linear mass density remains unchanged, the new wave speed is:}
\begin{choices}
\item $v/2$
\item $v$
\item $2v$
\item $4v$
\end{choices}

% QUESTION A6: Transmission coefficient
\mcq{A harmonic wave travels from string~1 (impedance $Z_1$) to string~2 (impedance $Z_2 = 3Z_1$). The amplitude transmission coefficient $t$ is:}
\begin{choices}
\item $1/2$
\item $2/3$
\item $3/4$
\item $3/2$
\end{choices}

% QUESTION A7: Sound intensity and decibels
\mcq{The threshold of hearing is $I_0 = 10^{-12}$~W/m$^2$. What is the intensity of a sound measured at 60~dB?}
\begin{choices}
\item $10^{-9}$~W/m$^2$
\item $10^{-8}$~W/m$^2$
\item $10^{-6}$~W/m$^2$
\item $10^{-3}$~W/m$^2$
\end{choices}

% QUESTION A8: Doppler with reflection off moving wall
\mcq{A stationary source emits sound of frequency $f_0 = 500$~Hz. The sound reflects off a wall that is moving towards the source at speed $v_w = 10$~m/s. Taking the speed of sound in air to be $v = 340$~m/s, the frequency of the reflected sound heard by the source is closest to:}
\begin{choices}
\item $500$~Hz
\item $515$~Hz
\item $530$~Hz
\item $545$~Hz
\end{choices}

% QUESTION A9: Pipe harmonics comparison
\mcq{An open-open pipe and an open-closed pipe have the same length $L$. The speed of sound is $v$. The ratio of the second harmonic frequency of the open-open pipe to the fundamental frequency of the open-closed pipe is:}
\begin{choices}
\item $2$
\item $4$
\item $8$
\item $1$
\end{choices}

% QUESTION A10: Beats
\mcq{Two identical guitar strings are plucked simultaneously. String~A is perfectly tuned to 440~Hz. String~B has a slightly higher tension than string~A. The musician hears beats at a frequency of 3~Hz. Which of the following is true?}
\begin{choices}
\item String~B has frequency 437~Hz
\item String~B has frequency 443~Hz
\item String~B could have frequency 437~Hz or 443~Hz
\item The beat frequency depends on the amplitudes of the two strings
\end{choices}

\newpage

% ============================================================
% SECTION B: QUANTUM PHYSICS MCQ
% ============================================================
\section*{Section B: Quantum Physics}
\noindent\textit{Answer all 10 questions. Each question is worth 2 marks.}\\[0.3cm]

\setcounter{mcqcounter}{0}

% QUESTION B1: Ket notation / outer product
\mcq{In Dirac notation, the object $\ket{\psi}\bra{\phi}$ is:}
\begin{choices}
\item A scalar (complex number)
\item A ket (state vector)
\item An operator
\item A bra (dual vector)
\end{choices}

% QUESTION B2: Global phase
\mcq{Let $\ket{\psi}$ be a normalised quantum state and $\theta$ a real number. Are the states $e^{i\theta}\ket{\psi}$ and $\ket{\psi}$ physically distinguishable?}
\begin{choices}
\item Yes, they have different norms
\item Yes, they give different measurement probabilities
\item No, they differ only by a global phase and are physically equivalent
\item Only if $\theta = \pi$
\end{choices}

% QUESTION B3: Trace
\mcq{For two normalised states $\ket{a}$ and $\ket{b}$ in a finite-dimensional Hilbert space, what is $\text{Tr}\!\left(\ketbra{a}{b}\right)$?}
\begin{choices}
\item $1$
\item $0$
\item $\braket{b}{a}$
\item $\braket{a}{b}$
\end{choices}

% QUESTION B4: Commutator of spin operators
\mcq{For a spin-$\tfrac{1}{2}$ particle, the commutator $[S_x, S_y]$ is:}
\begin{choices}
\item $0$
\item $i\hbar S_z$
\item $\tfrac{i\hbar}{2}\,S_z$
\item $2i\hbar\,S_z$
\end{choices}

% QUESTION B5: Probability from state
\mcq{A spin-$\tfrac{1}{2}$ particle is in the state $\ket{\psi} = \dfrac{1}{\sqrt{3}}\ket{\!\uparrow_z} + \sqrt{\dfrac{2}{3}}\,\ket{\!\downarrow_z}$. The probability of measuring $S_z$ and obtaining $-\hbar/2$ is:}
\begin{choices}
\item $1/3$
\item $1/\sqrt{3}$
\item $2/3$
\item $1/2$
\end{choices}

% QUESTION B6: Sequential measurements
\mcq{A measurement of $S_z$ on a spin-$\tfrac{1}{2}$ particle yields the result $+\hbar/2$, so the post-measurement state is $\ket{\!\uparrow_z}$. A subsequent measurement of $S_y$ is immediately performed. The probability of obtaining $+\hbar/2$ for this $S_y$ measurement is:}
\begin{choices}
\item $0$
\item $1/4$
\item $1/2$
\item $1$
\end{choices}

% QUESTION B7: Stationary vs non-stationary states
\mcq{Which of the following statements about stationary states is correct?}
\begin{choices}
\item A stationary state is any normalised quantum state
\item A stationary state is an eigenstate of the Hamiltonian; measurement probabilities in \emph{any} basis are time-independent
\item A stationary state can be a superposition of energy eigenstates with different energies
\item A stationary state has time-independent probabilities only when measured in the energy eigenbasis
\end{choices}

% QUESTION B8: Tensor product dimension
\mcq{Two spin-$\tfrac{1}{2}$ particles form a composite system. The dimension of the composite state space is:}
\begin{choices}
\item $2$
\item $3$
\item $4$
\item $8$
\end{choices}

% QUESTION B9: Pauli exclusion principle
\mcq{Two electrons cannot occupy the same quantum state simultaneously. This is a consequence of the fact that electrons are:}
\begin{choices}
\item Bosons, which must be in symmetric states
\item Fermions, which must be in antisymmetric states under exchange
\item Charged particles that repel each other
\item Particles with spin~$\tfrac{1}{2}$ that obey the uncertainty principle
\end{choices}

% QUESTION B10: Beta decay
\mcq{In beta-minus decay, a neutron decays as $n \to p + e^- + \bar{\nu}_e$. Which fundamental force mediates this process?}
\begin{choices}
\item The electromagnetic force
\item The strong nuclear force
\item The weak nuclear force
\item Gravity
\end{choices}

\newpage

% ============================================================
% SECTION C: LONG ANSWER QUESTIONS
% ============================================================
\section*{Section C: Long Answer Questions}
\noindent\textit{Answer all 4 questions. Show all working.}

\setcounter{questioncounter}{0}

% ---- QUESTION C1: MECHANICS (15 marks) ----
\question{15}

\noindent\textbf{Newton's Laws, Friction, and Circular Motion}

\begin{parts}
\item A block of mass $m_1 = 5$~kg sits on a horizontal surface. The coefficient of static friction is $\mu_s = 0.4$ and the coefficient of kinetic friction is $\mu_k = 0.3$. A horizontal force $F$ is applied to the block.
\begin{parts}
\item What is the maximum horizontal force that can be applied before the block begins to move? \qmarks{2}
\item If a horizontal force $F = 25$~N is applied, find the acceleration of the block. \qmarks{2}
\end{parts}

\item The same 5~kg block is now placed on a \emph{frictionless} inclined plane at angle $\theta = 30^\circ$ to the horizontal. It is connected by a light, inextensible string passing over a frictionless pulley at the top of the incline to a hanging mass $m_2 = 3$~kg. Find the acceleration of the system and the tension in the string. \qmarks{4}

\item A car of mass $M = 1200$~kg travels at constant speed around a banked curve of radius $r = 80$~m. The banking angle is $\alpha = 20^\circ$ and no friction is required. Find the speed of the car (the ``design speed''). \qmarks{3}

\item A conical pendulum consists of a mass $m = 0.5$~kg on a string of length $L = 1.0$~m. The mass revolves in a horizontal circle, and the string makes a constant angle $\theta = 25^\circ$ with the vertical.
Find the tension in the string, the radius of the circular orbit, and the period of revolution. \qmarks{4}
\end{parts}

\newpage

% ---- QUESTION C2: PROPERTIES OF MATTER (15 marks) ----
\question{15}

\noindent\textbf{Maxwell--Boltzmann, Lennard-Jones, and Crystal Structure}

\begin{parts}
\item The Maxwell--Boltzmann speed distribution gives three characteristic speeds for molecules of mass $m$ at temperature $T$. State the formulae for: the most probable speed $c_\mathrm{mp}$, the mean speed $c_\mathrm{avg}$, and the root-mean-square speed $c_\mathrm{rms}$. (These are on the formula sheet.) \qmarks{3}

\item Calculate all three speeds for nitrogen molecules (N$_2$, molar mass $M = 28$~g/mol) at $T = 300$~K. Verify that $c_\mathrm{mp} < c_\mathrm{avg} < c_\mathrm{rms}$. \qmarks{4}

\item The Lennard-Jones potential is given by
\[
V(r) = 4\varepsilon\left[\left(\frac{a}{r}\right)^{12} - \left(\frac{a}{r}\right)^{6}\right].
\]
By differentiating and setting $\mathrm{d}V/\mathrm{d}r = 0$, find the equilibrium separation $r_\mathrm{eq}$ in terms of $a$. Hence find the depth of the potential well (the binding energy $|V(r_\mathrm{eq})|$) in terms of $\varepsilon$. \qmarks{4}

\item Iron has a body-centred cubic (BCC) crystal structure with lattice parameter $a_\mathrm{latt} = 2.87 \times 10^{-10}$~m. The atomic mass of iron is $55.85$~u, where $1\;\text{u} = 1.661 \times 10^{-27}$~kg. A BCC unit cell contains 2 atoms. Calculate the density of iron. \qmarks{4}
\end{parts}

\newpage

% ---- QUESTION C3: OSCILLATIONS AND WAVES (15 marks) ----
\question{15}

\noindent\textbf{SHM, Superposition, and Organ Pipes}

\begin{parts}
\item A block of mass $m = 0.4$~kg is attached to a horizontal spring with spring constant $k = 160$~N/m on a frictionless surface. The block is pulled 8~cm from equilibrium and released from rest. Write the equation $x(t)$ for the subsequent motion. Find the maximum velocity and the maximum acceleration. \qmarks{3}

\item The same block--spring system is now given an initial displacement $x_0 = 6$~cm and an initial velocity $v_0 = 120$~cm/s in the $+x$ direction at $t = 0$. Find the amplitude $A$ and phase constant $\phi$ of the resulting oscillation, and write the equation $x(t)$. \qmarks{4}

\item Two harmonic waves on the same string have the same amplitude $A = 3$~mm and frequency $f = 200$~Hz, but differ in phase by $\delta = \pi/3$. Both waves travel in the $+x$ direction with speed $v = 100$~m/s. Using phasor addition (or the superposition formula), find the resultant amplitude and write the equation for the combined wave. \qmarks{4}

\item A pipe of length $L = 1.7$~m is open at both ends. The speed of sound is $v = 340$~m/s.
\begin{parts}
\item Find the fundamental frequency and the first three harmonics (i.e.\ $f_1$, $f_2$, $f_3$, $f_4$). \qmarks{2}
\item The pipe is now closed at one end. Find the new fundamental frequency and the first three resonant frequencies. \qmarks{2}
\end{parts}
\end{parts}

\newpage

% ---- QUESTION C4: QUANTUM PHYSICS (15 marks) ----
\question{15}

\noindent\textbf{Time Evolution and Multi-Particle States}

\begin{parts}
\item A spin-$\tfrac{1}{2}$ particle has Hamiltonian
\[
H = \omega_0\,S_z = \frac{\omega_0\hbar}{2}\left(\ketbra{\!\uparrow}{\uparrow\!} - \ketbra{\!\downarrow}{\downarrow\!}\right).
\]
At $t = 0$ the particle is prepared in the state
\[
\ket{\psi(0)} = \cos\alpha\,\ket{\!\uparrow} + \sin\alpha\,\ket{\!\downarrow}, \qquad \alpha = \frac{\pi}{6}.
\]
Find $\ket{\psi(t)}$. Calculate the probabilities $P(\uparrow_z,\,t)$ and $P(\downarrow_z,\,t)$ and comment on their time dependence. \qmarks{4}

\item For the state in part~(a), calculate the expectation value $\expect{S_x}(t)$. At what times is $\expect{S_x} = 0$? \qmarks{4}

\item Consider two \emph{distinguishable} spin-$\tfrac{1}{2}$ particles. Write down a complete orthonormal basis for the composite system. How many basis states are there? Write down one example of an entangled state in this basis. \qmarks{3}

\item Two identical bosons can each be in one of three single-particle states $\ket{a}$, $\ket{b}$, $\ket{c}$. List all possible two-boson states. How many are there? Repeat for two identical fermions. \qmarks{4}
\end{parts}

\newpage

% ============================================================
% FORMULA SHEET
% ============================================================
% EQUATION SHEET — PHYS10012 Core Physics I
% This file is \input{} at the end of each exam paper

\newpage
\section*{Formula Sheet}
\noindent\rule{\textwidth}{0.4pt}

\subsection*{Mathematical Formulae}

\begin{align*}
&\text{Binomial expansion:} & (1+x)^n &\approx 1 + nx + \frac{n(n-1)}{2!}x^2 + \cdots \quad (|x| \ll 1) \\[4pt]
&\text{Euler's formula:} & e^{i\theta} &= \cos\theta + i\sin\theta \\[4pt]
&\text{Trig identities:} & \sin A + \sin B &= 2\cos\!\left(\tfrac{A-B}{2}\right)\sin\!\left(\tfrac{A+B}{2}\right) \\
& & \cos A + \cos B &= 2\cos\!\left(\tfrac{A-B}{2}\right)\cos\!\left(\tfrac{A+B}{2}\right)
\end{align*}

\subsection*{Mechanics}

\begin{align*}
&\text{SUVAT:} & v &= u + at, \quad s = ut + \tfrac{1}{2}at^2, \quad v^2 = u^2 + 2as, \quad s = \tfrac{1}{2}(u+v)t \\[4pt]
&\text{Dot product:} & \mathbf{a}\cdot\mathbf{b} &= |\mathbf{a}||\mathbf{b}|\cos\theta = a_x b_x + a_y b_y + a_z b_z \\[4pt]
&\text{Cross product:} & \mathbf{a}\times\mathbf{b} &= |\mathbf{a}||\mathbf{b}|\sin\theta\;\hat{\mathbf{n}} = \begin{vmatrix}\hat{\mathbf{i}}&\hat{\mathbf{j}}&\hat{\mathbf{k}}\\a_x&a_y&a_z\\b_x&b_y&b_z\end{vmatrix} \\[4pt]
&\text{Rocket equation:} & \Delta v &= u \ln\!\left(\frac{m_0}{m_f}\right) \\[4pt]
&\text{Gravitational PE:} & U &= -\frac{GMm}{r} \\[4pt]
&\text{Escape velocity:} & v_e &= \sqrt{\frac{2GM}{R}} \\[4pt]
&\text{Force from PE:} & F &= -\frac{dU}{dx} \\[4pt]
&\text{Elastic collision (1D):} & v_1' &= \frac{m_1 - m_2}{m_1 + m_2}v_1 + \frac{2m_2}{m_1+m_2}v_2 \\[4pt]
&\text{Centre of mass:} & \mathbf{R}_\text{cm} &= \frac{\sum m_i\mathbf{r}_i}{\sum m_i} \\[4pt]
&\text{Reduced mass:} & \mu &= \frac{m_1 m_2}{m_1 + m_2} \\[4pt]
&\text{Centripetal accel:} & a_c &= \frac{v^2}{r} = \omega^2 r \\[4pt]
&\text{Parallel axis thm:} & I &= I_\text{cm} + Md^2 \\[4pt]
&\text{Atwood machine:} & a &= \frac{(m_1-m_2)g}{m_1+m_2}, \quad T = \frac{2m_1 m_2 g}{m_1+m_2}
\end{align*}

\noindent\textbf{Moments of inertia} (about axis through centre of mass):
\begin{center}
\begin{tabular}{ll}
Thin ring (radius $R$): & $I = MR^2$ \\
Solid disk (radius $R$): & $I = \tfrac{1}{2}MR^2$ \\
Solid sphere (radius $R$): & $I = \tfrac{2}{5}MR^2$ \\
Hollow sphere (radius $R$): & $I = \tfrac{2}{3}MR^2$ \\
Thin rod (length $L$, about centre): & $I = \tfrac{1}{12}ML^2$ \\
Thin rod (length $L$, about end): & $I = \tfrac{1}{3}ML^2$ \\
\end{tabular}
\end{center}

\subsection*{Properties of Matter}

\begin{align*}
&\text{Ideal gas:} & PV &= nRT \\[4pt]
&\text{Heat capacities:} & C_P - C_V &= R \quad\text{(per mole)} \\[4pt]
&\text{Adiabatic:} & PV^\gamma &= \text{const}, \quad TV^{\gamma-1} = \text{const} \\[4pt]
&\text{Isothermal work:} & W &= nRT\ln\!\left(\frac{V_2}{V_1}\right) \\[4pt]
&\text{Carnot efficiency:} & \eta &= 1 - \frac{T_c}{T_h} \\[4pt]
&\text{Entropy:} & dS &= \frac{dQ_\text{rev}}{T}, \quad \Delta S = mc\ln\!\left(\frac{T_2}{T_1}\right), \quad \Delta S = nR\ln\!\left(\frac{V_2}{V_1}\right) \\[4pt]
&\text{Lennard-Jones:} & V(r) &= 4\varepsilon\left[\left(\frac{a}{r}\right)^{12} - \left(\frac{a}{r}\right)^{6}\right] \\[4pt]
&\text{Mean KE:} & \langle KE \rangle &= \tfrac{3}{2}k_B T \\[4pt]
&\text{MB speeds:} & c_\text{mp} &= \sqrt{\frac{2k_BT}{m}}, \quad c_\text{avg} = \sqrt{\frac{8k_BT}{\pi m}}, \quad c_\text{rms} = \sqrt{\frac{3k_BT}{m}}
\end{align*}

\subsection*{Oscillations and Waves}

\begin{align*}
&\text{SHM:} & x(t) &= A\cos(\omega_0 t + \phi), \quad \omega_0 = \sqrt{k/m} \\[4pt]
&\text{Damped SHM:} & x(t) &= Ae^{-\gamma t/2}\cos(\omega_d t + \phi), \quad \omega_d = \sqrt{\omega_0^2 - \gamma^2/4} \\[4pt]
&\text{Quality factor:} & Q &= \omega_0/\gamma \\[4pt]
&\text{Forced amplitude:} & A(\omega) &= \frac{F_0/m}{\sqrt{(\omega_0^2 - \omega^2)^2 + \gamma^2\omega^2}} \\[4pt]
&\text{Phase lag:} & \tan\delta &= \frac{\gamma\omega}{\omega_0^2 - \omega^2} \\[4pt]
&\text{Pendulums:} & T_\text{simple} &= 2\pi\sqrt{L/g}, \quad T_\text{physical} = 2\pi\sqrt{I/(Mgh)} \\[4pt]
&\text{Wave equation:} & \frac{\partial^2 y}{\partial x^2} &= \frac{1}{v^2}\frac{\partial^2 y}{\partial t^2} \\[4pt]
&\text{String speed:} & v &= \sqrt{T/\mu} \\[4pt]
&\text{Sound in gas:} & v &= \sqrt{\gamma P/\rho} \\[4pt]
&\text{Harmonic wave:} & y &= A\sin(kx - \omega t) \\[4pt]
&\text{String impedance:} & Z &= \mu v = \sqrt{\mu T} \\[4pt]
&\text{Reflection:} & r &= \frac{Z_1 - Z_2}{Z_1 + Z_2}, \quad t = \frac{2Z_1}{Z_1+Z_2} \\[4pt]
&\text{Power (string):} & \langle P\rangle &= \tfrac{1}{2}\mu\omega^2 A^2 v \\[4pt]
&\text{Decibels:} & \beta &= 10\log_{10}(I/I_0), \quad I_0 = 10^{-12}\;\text{W/m}^2 \\[4pt]
&\text{Doppler:} & f' &= f\frac{v \pm v_o}{v \mp v_s} \\[4pt]
&\text{Group velocity:} & v_g &= \frac{d\omega}{dk} \\[4pt]
&\text{String harmonics:} & f_n &= \frac{nv}{2L} \quad (n=1,2,3,\ldots) \\[4pt]
&\text{Open-closed pipe:} & f_n &= \frac{nv}{4L} \quad (n=1,3,5,\ldots)
\end{align*}

\subsection*{Quantum Physics}

\begin{align*}
&\text{Schr\"{o}dinger eq:} & i\hbar\frac{d}{dt}\ket{\psi(t)} &= H\ket{\psi(t)} \\[4pt]
&\text{Time evolution:} & \ket{\psi(t)} &= \sum_n c_n e^{-iE_n t/\hbar}\ket{E_n}, \quad U(t) = e^{-iHt/\hbar} \\[4pt]
&\text{Spin operators:} & S_z &= \frac{\hbar}{2}\bigl(\ketbra{\!\uparrow}{\uparrow\!} - \ketbra{\!\downarrow}{\downarrow\!}\bigr) \\
& & S_x &= \frac{\hbar}{2}\bigl(\ketbra{\!\uparrow}{\downarrow\!} + \ketbra{\!\downarrow}{\uparrow\!}\bigr) \\
& & S_y &= \frac{\hbar}{2}\bigl(-i\ketbra{\!\uparrow}{\downarrow\!} + i\ketbra{\!\downarrow}{\uparrow\!}\bigr) \\[4pt]
&\text{Spin eigenstates:} & \ket{\!\uparrow_x} &= \tfrac{1}{\sqrt{2}}\bigl(\ket{\!\uparrow_z} + \ket{\!\downarrow_z}\bigr), \quad \ket{\!\downarrow_x} = \tfrac{1}{\sqrt{2}}\bigl(\ket{\!\uparrow_z} - \ket{\!\downarrow_z}\bigr) \\
& & \ket{\!\uparrow_y} &= \tfrac{1}{\sqrt{2}}\bigl(\ket{\!\uparrow_z} + i\ket{\!\downarrow_z}\bigr), \quad \ket{\!\downarrow_y} = \tfrac{1}{\sqrt{2}}\bigl(\ket{\!\uparrow_z} - i\ket{\!\downarrow_z}\bigr) \\[4pt]
&\text{Pauli matrices:} & \sigma_x &= \begin{pmatrix}0&1\\1&0\end{pmatrix}, \quad \sigma_y = \begin{pmatrix}0&-i\\i&0\end{pmatrix}, \quad \sigma_z = \begin{pmatrix}1&0\\0&-1\end{pmatrix}
\end{align*}

\subsection*{Physical Constants}

\begin{center}
\begin{tabular}{lll}
Speed of light & $c$ & $3.00 \times 10^{8}$ m/s \\
Planck's constant & $h$ & $6.626 \times 10^{-34}$ J$\cdot$s \\
Reduced Planck's constant & $\hbar$ & $1.055 \times 10^{-34}$ J$\cdot$s \\
Boltzmann constant & $k_B$ & $1.381 \times 10^{-23}$ J/K \\
Gas constant & $R$ & $8.314$ J/(mol$\cdot$K) \\
Avogadro's number & $N_A$ & $6.022 \times 10^{23}$ mol$^{-1}$ \\
Gravitational constant & $G$ & $6.674 \times 10^{-11}$ N$\cdot$m$^2$/kg$^2$ \\
Acceleration due to gravity & $g$ & $9.81$ m/s$^2$ \\
Electron mass & $m_e$ & $9.109 \times 10^{-31}$ kg \\
Proton mass & $m_p$ & $1.673 \times 10^{-27}$ kg \\
Elementary charge & $e$ & $1.602 \times 10^{-19}$ C \\
Mass of Earth & $M_\oplus$ & $5.97 \times 10^{24}$ kg \\
Radius of Earth & $R_\oplus$ & $6.37 \times 10^{6}$ m \\
\end{tabular}
\end{center}

\noindent\rule{\textwidth}{0.4pt}


\end{document}
